% GENERATED FILE. Edit canonical lessons, then run npm run generate:books.
% canonical-source-hash: fnv1a64:e878033d44b641d7
% canonical-lessons: TE-C58-a-little, TE-C58-more, TE-C58-less, TE-C58-no-thanks, TE-C58-thats-all

\chapter{More Short Replies}
\label{ch:te-more-replies}
\hlchaptermodality{58}

\begin{chapteropening}
\textbf{By the end of this chapter:} \emph{I can say how much I want, refuse politely, and close a conversation in Telugu.}

Complete TE-C58-thats-all and demonstrate the chapter promise.
\end{chapteropening}

\section[koñcem]{\te{కొంచెం} (koñcem) — a little, a bit}

\label{lesson:TE-C58-a-little}



\begin{quote}
Before the new one: what did \emph{emuka} mean?
\end{quote}

\subsection*{You'll want to know: \te{కొంచెం}}
\textbf{\te{కొంచెం}} (\emph{koñcem}) — \textquotedblleft{}a little, a bit\textquotedblright{}.

\te{కొంచెం} is \textquotedblleft{}a little\textquotedblright{}, and it does far more work in Telugu than its size suggests. It softens whatever comes after it.

Ask for something plainly and you are demanding; put \te{కొంచెం} in front and you are asking. \te{కొంచెం} \te{నీళ్ళు} is a request for water that leaves the other person room. The word is a quantity in the dictionary and a manner in the mouth.

Telugu shares this habit with a good many languages. English does it with \textquotedblleft{}just\textquotedblright{} and \textquotedblleft{}a bit\textquotedblright{}, and both of those have stopped measuring anything too.

\begin{grammarlens}[title={What you've built}]
A word that makes the next one polite.
\end{grammarlens}

\subsection*{Your turn}
Say these aloud:

\begin{itemize}
\raggedright
  \item \emph{koñcem}
  \item it once more, slowly
  \item \emph{koñcem}, then \emph{koñcem pālu}, and hear how much softer it is than \emph{pālu} alone
\end{itemize}

\begin{tcolorbox}[breakable,colback=teal!4,colframe=teal!35!black,title={Before you move on}]
What does \te{కొంచెం} mean? (\textquotedblleft{}a little, a bit\textquotedblright{}.) Say it once more.
\end{tcolorbox}
\section[ekkuva]{\te{ఎక్కువ} (ekkuva) — more, a great deal}

\label{lesson:TE-C58-more}



\begin{quote}
Before the new one: what did \emph{koñcem} mean?
\end{quote}

\subsection*{You'll want to know: \te{ఎక్కువ}}
\textbf{\te{ఎక్కువ}} (\emph{ekkuva}) — \textquotedblleft{}more, a great deal\textquotedblright{}.

\te{ఎక్కువ} is built on \te{ఎక్కు}, \textquotedblleft{}to climb, to go up\textquotedblright{}. The quantity is a movement: more is what has gone up.

Set it against \te{కొంచెం}, which you had a moment ago, and you have the two ends of a plate of food. \te{కొంచెం} is what you say when you would like the serving to stop; \te{ఎక్కువ} is what you say when it was too much.

It doubles as \textquotedblleft{}too much\textquotedblright{} without changing shape, so tone and face carry the difference. \te{ఎక్కువ} \te{మాట్లాడు} can be praise or a complaint depending on how it is said.

\begin{grammarlens}[title={What you've built}]
Two: a little, and a lot.
\end{grammarlens}

\subsection*{Your turn}
Say these aloud:

\begin{itemize}
\raggedright
  \item \emph{ekkuva}
  \item it once more, slowly
  \item \emph{ekkuva}, then \emph{koñcem}, and say them as a pair
\end{itemize}

\begin{tcolorbox}[breakable,colback=teal!4,colframe=teal!35!black,title={Before you move on}]
What does \te{ఎక్కువ} mean? (\textquotedblleft{}more, a great deal\textquotedblright{}.) Say it once more.
\end{tcolorbox}
\section[takkuva]{\te{తక్కువ} (takkuva) — less, too little}

\label{lesson:TE-C58-less}



\begin{quote}
Before the new one: what did \emph{ekkuva} mean?
\end{quote}

\subsection*{You'll want to know: \te{తక్కువ}}
\textbf{\te{తక్కువ}} (\emph{takkuva}) — \textquotedblleft{}less, too little\textquotedblright{}.

\te{తక్కువ} is built on \te{తక్కు}, \textquotedblleft{}to fall short\textquotedblright{}, exactly the way \te{ఎక్కువ} is built on climbing. Two words, one shape of construction, opposite directions.

A matched pair like this is worth holding on to, because it lets you guess. Once the -\te{కువ} ending on both has registered, you can meet a Telugu word for the first time and have some idea what it is doing.

\te{తక్కువ} \te{చేసి} \te{మాట్లాడు}, to speak somebody less, is to belittle them. The arithmetic goes straight into the manners.

\begin{grammarlens}[title={What you've built}]
Three, and two of them are a pair.
\end{grammarlens}

\subsection*{Your turn}
Say these aloud:

\begin{itemize}
\raggedright
  \item \emph{takkuva}
  \item it once more, slowly
  \item \emph{takkuva}, then \emph{ekkuva}, and say which one climbs
\end{itemize}

\begin{tcolorbox}[breakable,colback=teal!4,colframe=teal!35!black,title={Before you move on}]
What does \te{తక్కువ} mean? (\textquotedblleft{}less, too little\textquotedblright{}.) Say it once more.
\end{tcolorbox}
\section[vaddu]{\te{వద్దు} (vaddu) — no, I don't want it}

\label{lesson:TE-C58-no-thanks}



\begin{quote}
Before the new one: what did \emph{takkuva} mean?
\end{quote}

\subsection*{You'll want to know: \te{వద్దు}}
\textbf{\te{వద్దు}} (\emph{vaddu}) — \textquotedblleft{}no, I don't want it\textquotedblright{}.

\te{వద్దు} is the refusal of a thing offered. It is not the same as \te{కాదు}, which denies that something is so; \te{వద్దు} turns down an offer.

It answers \te{కావాలి}, \textquotedblleft{}I want it\textquotedblright{}, and the two of them run every meal in a Telugu house. A hand over the plate and \te{వద్దు} is a complete sentence.

It softens the moment you put \te{కొంచెం}'s cousin in front of it, or the \textbf{-\te{అండి}} ending on the end: \te{వద్దండి} is the same refusal with the door left open.

\begin{grammarlens}[title={What you've built}]
Four.
\end{grammarlens}

\subsection*{Your turn}
Say these aloud:

\begin{itemize}
\raggedright
  \item \emph{vaddu}
  \item it once more, slowly
  \item \emph{vaddu}, then \emph{vaddaṇḍi}, and say which one you would use with a stranger
\end{itemize}

\begin{tcolorbox}[breakable,colback=teal!4,colframe=teal!35!black,title={Before you move on}]
What does \te{వద్దు} mean? (\textquotedblleft{}no, I don't want it\textquotedblright{}.) Say it once more.
\end{tcolorbox}
\section[antē]{\te{అంతే} (antē) — that's all; that's just it}

\label{lesson:TE-C58-thats-all}



\begin{quote}
Before the new one: what did \emph{vaddu} mean?
\end{quote}

\subsection*{You'll want to know: \te{అంతే}}
\textbf{\te{అంతే}} (\emph{antē}) — \textquotedblleft{}that's all; that's just it\textquotedblright{}.

\te{అంతే} is \te{అంత}, \textquotedblleft{}that much\textquotedblright{}, with a small \textbf{-\te{ఏ}} on the end that means \textquotedblleft{}just, exactly\textquotedblright{}. So the word says \textquotedblleft{}exactly that much\textquotedblright{} — and it ends a list, a story or an argument.

It has a second life as agreement. Somebody explains why a thing went wrong and you say \te{అంతే}: that is precisely it, nothing further to add. The same word closes a sale and closes a discussion.

It is a good word to leave a chapter on. \te{కొంచెం}, \te{ఎక్కువ}, \te{తక్కువ}, \te{వద్దు} — and then \te{అంతే}.

\begin{grammarlens}[title={What you've built}]
Five: \te{కొంచెం}, \te{ఎక్కువ}, \te{తక్కువ}, \te{వద్దు}, \te{అంతే}. You can now say how much, and when to stop.
\end{grammarlens}

\subsection*{Your turn}
Say these aloud:

\begin{itemize}
\raggedright
  \item \emph{antē}
  \item it once more, slowly
  \item all five in order, then refuse a second helping with \emph{vaddu}, \emph{antē}
\end{itemize}

\begin{tcolorbox}[breakable,colback=teal!4,colframe=teal!35!black,title={Before you move on}]
What does \te{అంతే} mean? (\textquotedblleft{}that's all; that's just it\textquotedblright{}.) Say it once more.
\end{tcolorbox}
